\BeleskeID{09_3-s060}
% element: 09_3-s060:e05
% slide-id: 09_3-s060
Transmisioni gejt, odnosno bilateralni prekidač, koristi se radi izbegavanja degradacije naponskih nivoa, naročito u analognim delovima digitalnih sistema. Dodavanjem P tranzistora paralelno N tranzistoru obezbeđuje se prenos u različitim uslovima rada. Primena je naročito pogodna za selektorske funkcije.

Na simbolima nisu posebno označeni ulaz i izlaz: struktura je simetrična, pa signal može prolaziti u oba smera. Za razmatranje analognog prekidača kontrolni napon uzimamo kao $V_C=\pm15\,\mathrm V$, dok se ulaz menja između $-12\,\mathrm V$ i $+12\,\mathrm V$. Otpornik $R_L$ predstavlja opterećenje izlaza.
